\documentclass[11pt,a4paper]{article}

\usepackage{amsmath,amssymb,amsthm,mathtools,microtype,booktabs}
\usepackage[margin=25mm]{geometry}
\usepackage[hyphens]{url}
\usepackage[colorlinks=true,urlcolor=blue,linkcolor=blue,citecolor=blue]{hyperref}
\renewcommand{\baselinestretch}{1.08}
\newtheorem{theorem}{Theorem}
\newtheorem{lemma}{Lemma}
\newtheorem{corollary}{Corollary}
\theoremstyle{remark}
\newtheorem{remark}{Remark}
\newcommand{\R}{\mathbb R}
\newcommand{\N}{\mathbb N}
\newcommand{\Q}{\mathcal Q}
\numberwithin{equation}{section}
\title{Referee report on ``Uniform two-parameter asymptotics\\ for Lanczos
coefficients of centered uniform convolutions''\\[4pt]
\large with independent proofs, a partial Lean~4 formalisation and numerical checks}
\author{Independent review prepared for Roberto Alfano\\ by Claude (Anthropic), an AI system}
\date{22 October 2026}
\begin{document}
\maketitle

\noindent\textbf{Scope and status.} This is a separate research document. It reviews the
five-page note of 20 October 2026 (henceforth \emph{the note}) under the instructions of
the author. The manuscripts V032 and V035 were not modified. Mathematical validity,
bibliographic novelty and formal-machine certification are reported separately, in
Sections~\ref{sec:verdict}--\ref{sec:proof}, Section~\ref{sec:orig} and
Section~\ref{sec:lean} respectively. Statements of Levin and Lubinsky, Floater and
Milovanovi\'c were checked against the original documents supplied by the author
(page images were inspected wherever the extracted text was garbled). Numerical
computations are consistency checks only; they are not used as proof of anything.

\section{Referee's verdict.}\label{sec:verdict}

\textbf{Verdict: REPAIRABLE.} I believe the theorem is true, and the architecture of the
proof is sound: the scaling, the reduction to the Mhaskar--Rakhmanov--Saff equations and
the use of Levin--Lubinsky's Theorem~13.4 at the degrees $m=N-1,N$ are correct, and every
factor of $M$, $N$, $2$ and $\pi$ checks out (Section~\ref{sec:norm}). No counterexample
was found; on the contrary, high-precision computations up to $M=400$, $N=1600$ show
$\beta_N^{[M]}/M^2-F(N/M)=c(N/M)/M+O(M^{-2})$ with a stable function $c$
(Section~\ref{sec:num}). The note is not VALID as submitted because several steps are
asserted rather than proved. The most substantial one is the remainder estimate in the
differentiated saddlepoint expansion of Lemma~1, which is the analytic heart of the
argument. There are also three smaller defects. Lemma~2 is false as literally stated for
$M=1$. The choice of $r^*$ in Lemma~3 does not account for $a_{N,M}\neq a_t$. Definition~1.1
of Levin--Lubinsky requires a whole sequence indexed by every $n\ge1$, which the diagonal
argument must supply explicitly. None of these gaps is fatal. Section~\ref{sec:proof}
supplies a complete proof of each missing step; with these repairs the theorem is, in my
assessment, proved on paper, \emph{conditionally on Theorem~13.4 of Levin--Lubinsky as
published}. It is \emph{not} certified by machine in full: the Lean~4 development covers
the logical skeleton and every elementary step, but not the Fourier--Laplace analysis of
Lemma~1 nor the potential theory behind Levin--Lubinsky's theorem (Section~\ref{sec:lean}).

\section{Line-by-line mathematical audit.}\label{sec:audit}

Each finding gives the place in the note, the issue, the source checked and a severity:
\emph{fatal} (invalidates the result), \emph{substantial} (a step of the proof is missing),
\emph{minor} (a correction of statement or wording). No fatal problem was found.

\paragraph{A1. Page 1, definitions (minor, no change needed).} The density of $Y_M=X_M/M$ is
$g_M(x)=Mf_M(Mx)$. The quantity $\Lambda(s)=\log\{\sinh(s/2)/(s/2)\}$ is the cumulant
generating function of $\operatorname{Uniform}[-\frac12,\frac12]$, since
$\int_{-1/2}^{1/2}e^{st}\,dt=2\sinh(s/2)/s$. For the monic recurrence of the note,
$P_{N+1}=xP_N-\beta_NP_{N-1}$ (no diagonal term, by symmetry), the orthonormal recurrence
is $xp_N=b_{N+1}p_{N+1}+b_Np_{N-1}$ with $b_N=\beta_N^{1/2}$, which is the $N$th
off-diagonal entry of the Jacobi matrix. This is Milovanovi\'c's indexing
\cite[(2)--(3)]{milovanovic2017}, so the $\beta_k^{[m]}$ there and the $\beta_N^{[M]}$ of the
note coincide (his weight is $f_M$ translated to $[0,M]$, and translation does not change
$\beta$). In the theorem, $N\ge L$ and $N/M\in K$ force $M\ge L/\max K$, so both
parameters tend to infinity.

\paragraph{A2. Equation (3) of the note, Legendre duality (correct).} The identities $I'=s$
and $I''=1/v$ hold; they are formalised in Lean (\texttt{legendre\_deriv},
\texttt{legendre\_second}).

\paragraph{A3. Lemma~1, Fourier representation (6) (minor).} The formula is correct for
$M\ge2$. The cleanest justification needs no contour shift: the tilted law with density
$e^{sy}f_M(y)/A(s)^M$ has characteristic function $\{A(s+iu)/A(s)\}^M$, which is integrable
for $M\ge2$ because $|A(s+iu)|\le2\cosh(s/2)/|u|$, and Fourier inversion applies. The note
gives no justification; Section~\ref{sec:L1} supplies it.

\paragraph{A4. Lemma~1, analytic logarithm and (7) (minor).} The zeros of $A(z)=2\sinh(z/2)/z$
are $z=2\pi ik$, $k\ne0$, so $A$ does not vanish on a strip $\{\Re z\in S',\ |\Im z|\le
\varepsilon\}$ for every compact $S'$ and $\varepsilon<2\pi$; for small $\varepsilon$ the
values of $A$ there have positive real part and the principal logarithm is analytic. The
note asserts this without proof; it is true.

\paragraph{A5. Lemma~1, derivative cancellations (correct).} With $s'v=1$, the coefficient
of $iu$ in $\partial_x\Psi$ is $s'v-1=0$, and the coefficient of $iu$ in $\partial_x^2\Psi$
is $s''v+s'^2\Lambda'''(s)=0$. Both identities are proved in Lean
(\texttt{tilt\_first}, \texttt{tilt\_second}, \texttt{tilt\_cancellations}). Hence
$\partial_x\Psi=O_J(u^2)$ and $\partial_x^2\Psi=O_J(u^2)$, as claimed.

\paragraph{A6. Lemma~1, outer region (substantial, repairable).} The bounds
$\rho_{J,\varepsilon,R}<1$ and $|A(s+iu)/A(s)|\le C_J/(1+|u|)$ are correct. The claim that
$\partial_x^j$ of the integrand is ``bounded by a polynomial in $M$ and $1+|u|$ times an
$(M-j)$th power'' is correct but unproved. The required estimates are
$|\partial_xG|\le C$ and $|\partial_x^2G|\le C(1+|u|)$ for
$G=e^{-iux}A(s(x)+iu)/A(s(x))$; they are proved in Section~\ref{sec:L1}. The justification of
differentiating under the integral sign twice is not mentioned; it follows from the same
majorants for $M\ge7$.

\paragraph{A7. Lemma~1, inner region (substantial, repairable).} The expansion of
$e^{M\Psi}$ is correct, including the coefficients. The bound
$C_JM^{-3/2}(1+|z|^{12})e^{c_Jz^2/M^{1/3}}$ for $\mathcal E_M$ and its two $x$-derivatives
is stated without proof, and the exponential factor is unexplained (it is $O(1)$ on
$|z|\le M^{1/12}$ and unnecessary). This is the decisive estimate of the lemma, so its
absence is a real gap. Section~\ref{sec:L1} proves the bound $C(1+|z|)^{18}M^{-3/2}$.
The vanishing of the $M^{-1/2}$ term, also after $x$-differentiation, is correct because
the truncation $|z|\le M^{1/12}$ is symmetric. The conclusion is moreover confirmed to
second order by computation: the $C^2$ remainder of (4) in the note equals $q(x)/M+O(M^{-2})$ with the classical
correction $q=\kappa_4/(8v^2)-5\kappa_3^2/(24v^3)$, and $M$ times the remainder and its first
two derivatives converge (Table~\ref{tab:saddle}).

\paragraph{A8. Lemma~1, consequences (5) (correct).} With
$\log g_M=-MI+\frac12\log\{M/(2\pi v)\}+O_{C^2}(M^{-1})$ one gets
$\|\phi_M-I\|_{C^2(J)}=O(M^{-1}\log M)$ and $\|\phi_M'-I'\|_{C^1(J)}=O(M^{-1})$.

\paragraph{A9. Lemma~2, the limit function $H$ (correct, one step brief).} Continuity,
$H(0^+)=0$, strict increase and $H(a)\to\infty$ as $a\uparrow\frac12$ are correct. The
divergence deserves one line: $\Lambda'(s)=\frac12\coth(s/2)-1/s$ gives
$I'(x)\sim(\frac12-x)^{-1}$, and $\frac12-\frac12\sin\theta\sim\frac14(\frac\pi2-\theta)^2$
makes $\int^{\pi/2}d\theta/(1-\sin\theta)$ diverge. Strict increase and continuity are
formalised (\texttt{Hint\_strictMonoOn}, \texttt{Hint\_continuousOn}).

\paragraph{A10. Lemma~2, statement for small $M$ (minor).} As literally stated (``for every
$K$ and $t=N/M\in K$ there exists a unique $a_{N,M}$'') the lemma includes $M=1$, where
$g_1\equiv1$ on its support, $\phi_1'\equiv0$ and no positive root exists. For $M=2$ a
unique root exists, since $x\phi_2'(x)=|x|/(1-2|x|)$ is continuous and increasing, but $Q'$
is then discontinuous at $0$, so Lemma~3 needs $M\ge3$. The statement should read
``for every $M\ge2$'', and the estimate (9) of the note holds for $M\ge M_0(K)$.

\paragraph{A11. Lemma~2, unimodality and strict log-concavity (correct; reference
clarified).} The induction $f_{M+1}'(x)=f_M(x+\frac12)-f_M(x-\frac12)<0$ for
$0<x<(M+1)/2$ is correct with base case $M=2$ (not $M=1$). The argument that $\log f_M$ is
not affine on any subinterval is correct, because a polynomial of the form
$e^{\lambda x+\mu}$ on an interval is constant, and $f_M$ has exact degree $M-1\ge1$ on every
knot interval (the leading coefficient on the $(k+1)$th interval is
$(-1)^k\binom{M-1}{k}/(M-1)!\neq0$). Floater \cite[Lemma~2, (11) and the discussion after
it]{floater2024} gives more: for degree $d=M-1\ge2$ the B-spline is $d$-concave, inequality
(11) is strict on the interior knot intervals, and on the two end intervals the spline is a
pure power. Since $ff''-(f')^2=\{ff''-\frac{d-1}{d}(f')^2\}-\frac1d(f')^2$, this gives
$(\log f_M)''<0$ on every open knot interval, which is exactly what the note needs. It
makes $\phi_M'$ continuous and strictly increasing on the open support for $M\ge3$, where
$f_M\in C^1$.

\paragraph{A12. Lemma~2, estimate (9) (substantial for the write-up, now formalised).} The
note's proof omits the a priori localisation of $a_{N,M}$ in a compact subinterval where
$H'$ is bounded below. The complete argument, with explicit constants, is formalised in
Lean (\texttt{endpoint\_lemma}): if $|\phi_M'-I'|\le C/M$ on $[0,b]$ and
$H(\mathrm{lo})+bC/(2M)<t_1\le t_2<H(\mathrm{hi})-bC/(2M)$, then for every $t\in[t_1,t_2]$
the finite-$M$ root exists, is unique, lies in $(\mathrm{lo},\mathrm{hi})$ and satisfies
$|a_{N,M}-a_t|\le2bC/(m\,\mathrm{lo}\,M)$, where $m$ is a lower bound for $I''$ on $[0,b]$.

\paragraph{A13. Lemma~3, conditions (i)--(iv) of Definition~1.1 (correct, two
corrections).} I checked Definition~1.1 \cite[pp.~6--7]{levinlubinsky2018} directly.
Condition (i) holds because $Q_{N,M}\to+\infty$ at the finite endpoints $\pm1/(2a)$.
Condition (ii) requires $M\ge3$. Conditions (iii) and (iv) follow from a uniform bound on
$\phi_M''$ over $[-ar^*,ar^*]$ (Lemma~1) by the chain rule and Lipschitz~$\Rightarrow$~H\"older
(both formalised: \texttt{Qprime\_lipschitz}, \texttt{holder\_of\_lipschitz}). Two
corrections are needed. First, $r^*$ must satisfy $r^*\sup a_{N,M}<\frac12$, not merely
$r^*a_+<\frac12$; take $r^*(a_++\delta)<\frac12$ and $M$ so large that
$a_{N,M}\le a_++\delta$ (Lemma~2). Second, Definition~1.1 demands the conditions for
\emph{every} $n\ge1$ with the same constants, so the finitely many small $n$ and the degrees
$n$ that do not occur must be filled with admissible weights; see A16.

\paragraph{A14. Lemma~3, condition (v) and the MRS equations (correct).} Remark~(b) after
Definition~1.1 \cite[p.~6]{levinlubinsky2018} states that the support condition is
equivalent to (1.13)--(1.14). With
$e^{-2NQ}=ag_M(ay)$ one has $Q'(y)=\frac{aM}{2N}\phi_M'(ay)$ (formalised:
\texttt{Q\_hasDerivAt}). I verified, and formalised in the arcsine form
(\texttt{cos\_odd\_eq\_two\_sin}, \texttt{mrs\_first\_iff}), that the first MRS equation is
equivalent to the endpoint equation (8) of the note, and that the second holds by parity
(\texttt{mrs\_second}). The normalisation is therefore correct. The change of variable
$y=\cos\theta$ in the improper integrals of (1.13)--(1.14) is elementary but not formalised.

\paragraph{A15. Proof of the theorem, application of Theorem~13.4 (correct; a misprint in
the source).} I read Theorems~13.1--13.4 \cite[pp.~123--124]{levinlubinsky2018} and the
proof of Theorem~13.4 \cite[p.~138]{levinlubinsky2018}. The recurrence there is
$xp_{n,m}=A_{n,m}p_{n,m+1}+B_{n,m}p_{n,m}+A_{n,m-1}p_{n,m-1}$ with
$A_{n,m}=\gamma_{n,m}/\gamma_{n,m+1}$, and (13.8) holds uniformly for $|m-n|\le
n^{1/3-\varepsilon}$, the window (13.3). The coefficient linking degrees $N-1$ and $N$ is
therefore $A_{N,N-1}$, as in the note, and $m=N-1$ lies in the window. Theorem~2.3(c)
\cite[p.~12]{levinlubinsky2018} prints the recurrence with a misprint (the term
$A_{n,m}p_{n,m-1}$ should read $A_{n,m}p_{n,m+1}$); the note's reading agrees with Chapter~13. The scaling
(10) of the note, $\beta_N^{[M]}/M^2=a_{N,M}^2A_{N,N-1}^2$, is verified in
Section~\ref{sec:norm} and its algebra is formalised (\texttt{beta\_scaling}).

\paragraph{A16. Proof of the theorem, uniformity (correct in substance, details
missing).} The contradiction argument is valid. The ``filling'' needs two explicit
facts: (a) for every large $n$ there is $M$ with $n/M$ in a fixed interval
$[t_1,t_2]\subset(0,\infty)$, and (b) the remaining small $n$ can be given a fixed admissible
weight, for instance $Q_n(x)=x^2$ on $\R$, for which (1.13) reads
$\frac1\pi\int_{-1}^1 2x^2(1-x^2)^{-1/2}dx=1$. Both the abstract argument and fact (a) are
formalised (\texttt{uniform\_of\_sequential}, \texttt{fill\_ratio}).

\paragraph{A17. Corollary (correct).} $\Lambda(s)=s^2/24-s^4/2880+O(s^6)$,
$I'(x)=12x+\frac{144}{5}x^3+O(x^5)$, $H(a)=3a^2+\frac{27}{5}a^4+O(a^6)$ and
$F(\alpha)=\frac{\alpha}{12}-\frac{\alpha^2}{20}+O(\alpha^3)$ are all correct; the last two
steps are formalised (\texttt{Hint\_expansion}, \texttt{inversion}). Monotonicity,
$0<F<\frac1{16}$ and $F(\alpha)\to\frac1{16}$ follow from Lemma~2. The note calls these
``checks at the two ends of the scaling family''. They are properties of $F$ only; the
theorem does not imply the fixed-$N$ or fixed-$M$ limits, because it is uniform only on
compact ratio sets.

\paragraph{A18. Novelty section (correct, one understatement).} The description of the
fixed-$M$ theory and of Milovanovi\'c's conjectures is accurate. In fact the Taylor
coefficients of $F$ agree with the top-degree parts of \emph{all five} coefficients in
Conjecture~2, not only the first two (Section~\ref{sec:orig}).

\section{Independent proof and necessary repairs.}\label{sec:proof}

Throughout, $J=[-b,b]$ with $0<b<\frac12$, $s(x)$ is the tilt with $\Lambda'(s(x))=x$,
$v=\Lambda''\circ s$, $\kappa_j=\Lambda^{(j)}\circ s$, and constants $C,c$ depend only on $b$
(and later on $K$).

\subsection{Normalisations.}\label{sec:norm}

The weight of $y=Y_M/a$ with $a=a_{N,M}$ is $ag_M(ay)$ on $|y|<1/(2a)$, and $X_M=May$. If
$P_n^{y}$ are the monic orthogonal polynomials of $y$ with $h_n^{y}=\|P_n^{y}\|^2$, then
$P_n^{X}(X)=(Ma)^nP_n^{y}\{X/(Ma)\}$ are those of $X_M$ and $h_n^{X}=(Ma)^{2n}h_n^{y}$. Since
$\beta_N=h_N/h_{N-1}$ for a monic recurrence, and the orthonormal leading coefficients are
$\gamma_n=(h_n^{y})^{-1/2}$ with $A_{N,N-1}=\gamma_{N,N-1}/\gamma_{N,N}$ in the convention of
\cite[p.~124 and p.~138]{levinlubinsky2018},
\begin{equation}\label{eq:scaling}
 \beta_N^{[M]}=\frac{h_N^{X}}{h_{N-1}^{X}}=(Ma)^2\,\frac{h_N^{y}}{h_{N-1}^{y}}
 =(Ma)^2A_{N,N-1}^2 .
\end{equation}
Multiplying the weight by a constant (Levin--Lubinsky's normalisation $Q_n(t_n)=0$,
Remark~(f) after Definition~1.1) changes neither $\beta$ nor $A$. The external field is
\begin{equation}\label{eq:Q}
 Q_{N,M}(y)=\frac{M}{2N}\,\phi_M(ay)-\frac{\log a}{2N},\qquad
 Q_{N,M}'(y)=\frac{aM}{2N}\,\phi_M'(ay),\qquad \phi_M=-\frac1M\log g_M ,
\end{equation}
and, with $y=\cos\theta$ and the oddness of $\phi_M'$, the first MRS equation becomes
\begin{equation}\label{eq:mrs}
 \frac1\pi\int_0^\pi\cos\theta\,Q_{N,M}'(\cos\theta)\,d\theta
 =\frac MN\cdot\frac1\pi\int_0^{\pi/2}a\sin\theta\,\phi_M'(a\sin\theta)\,d\theta=1 ,
\end{equation}
which is the note's endpoint equation with $t=N/M$. Equations
\eqref{eq:scaling}--\eqref{eq:mrs} are formalised (Section~\ref{sec:lean}). The limit
$A_{N,N-1}\to\frac12$ is the Chebyshev value for $[-1,1]$, so
$\beta_N^{[M]}/M^2\to a^2/4$; as a sanity check, $\alpha\to0$ gives $F(\alpha)\sim\alpha/12$,
i.e.\ $\beta_N\approx NM/12=N\operatorname{Var}X_M$ (Hermite), and $\alpha\to\infty$ gives
$\frac1{16}$, i.e.\ $\beta_N\approx M^2/16=(\text{length}/4)^2$.

\subsection{Lemma 1: a complete proof.}\label{sec:L1}

\begin{lemma}\label{lem:L1}
For every $b<\frac12$ there are $M_0$ and $C$ such that for $M\ge M_0$
\begin{equation}\label{eq:L1}
 \Bigl\|\log g_M+MI-\frac12\log\frac{M}{2\pi v}\Bigr\|_{C^2([-b,b])}\le\frac CM .
\end{equation}
\end{lemma}

\begin{proof}
\emph{Step 0 (representation).} Let $A(z)=\int_{-1/2}^{1/2}e^{zt}dt=2\sinh(z/2)/z$, so that
$\int e^{zy}f_M(y)\,dy=A(z)^M$ and $A=e^\Lambda$ on $\R$. For real $s$ the probability density
$e^{sy}f_M(y)A(s)^{-M}$ has characteristic function $\{A(s+iu)/A(s)\}^M$. Since
$|2\sinh\{(s+iu)/2\}|\le2\cosh(s/2)$ and $|s+iu|\ge|u|$,
\begin{equation}\label{eq:Adecay}
 |A(s+iu)|\le\frac{2\cosh(s/2)}{|u|},
\end{equation}
so for $M\ge2$ the characteristic function is integrable, $f_M$ is continuous, and Fourier
inversion gives
$e^{sy}f_M(y)A(s)^{-M}=\frac1{2\pi}\int_\R\{A(s+iu)/A(s)\}^Me^{-iuy}du$. Taking $y=Mx$,
$s=s(x)$ and using $A(s)^Me^{-sMx}=e^{M\{\Lambda(s)-sx\}}=e^{-MI(x)}$,
\begin{equation}\label{eq:rep}
 g_M(x)=\frac{M}{2\pi}\,e^{-MI(x)}J_M(x),\qquad J_M(x)=\int_\R G(x,u)^M\,du,\qquad
 G(x,u)=e^{-iux}\,\frac{A(s(x)+iu)}{A(s(x))}.
\end{equation}

\emph{Step 1 (bounds on $G$).} The map $x\mapsto s(x)$ is real-analytic on
$(-\frac12,\frac12)$ (inverse function theorem, $\Lambda''>0$), and $S=s(J)$ is compact.
(a) $|G|\le1$, and $|G(x,u)|<1$ for $u\ne0$ because $G(x,\cdot)$ is the characteristic
function of a law with a density; by continuity and compactness, and by
\eqref{eq:Adecay}, for every $\varepsilon>0$ there is $\rho_\varepsilon<1$ with
$|G|\le\rho_\varepsilon$ for $x\in J$, $|u|\ge\varepsilon$, and $|G|\le C_1/|u|$ for
$|u|\ge1$. (b) Write
\begin{equation}\label{eq:dG}
 \partial_xG=-iuG+s'(x)\,e^{-iux}\,\frac{A'(s+iu)A(s)-A(s+iu)A'(s)}{A(s)^2}.
\end{equation}
Since $A'(z)=\cosh(z/2)/z-2\sinh(z/2)/z^2$, we have $|A'(s+iu)|\le C/|u|$ for $|u|\ge1$, and
$|uG|\le C$; hence $|\partial_xG|\le C$. Differentiating \eqref{eq:dG} once more, the only
term not bounded by a constant is $-iu\,\partial_xG$, so $|\partial_x^2G|\le C(1+|u|)$.

\emph{Step 2 (outer region).} For $j\le2$, $\partial_x^j(G^M)$ is a combination of
$MG^{M-1}\partial_x^jG$ and $M(M-1)G^{M-2}(\partial_xG)^2$, so
$|\partial_x^j(G^M)|\le CM^2(1+|u|)|G|^{M-2}$. For $|u|\ge\varepsilon$ and $M\ge7$,
$|G|^{M-2}\le\rho_\varepsilon^{M-6}|G|^4$ and $(1+|u|)|G|^4$ is integrable by (a). This
majorant, uniform in $x\in J$, justifies differentiating $J_M$ twice under the integral and
gives
\begin{equation}\label{eq:outer}
 \sup_{x\in J}\,\Bigl|\partial_x^j\int_{|u|\ge\varepsilon}G^M\,du\Bigr|\le CM^2\rho_\varepsilon^{M-6},
 \qquad j=0,1,2 .
\end{equation}

\emph{Step 3 (inner region).} Choose $\varepsilon\in(0,1]$ so small that $A(\sigma+iu)$ has
positive real part for $\sigma$ in a neighbourhood $S'$ of $S$ and $|u|\le\varepsilon$
(possible because $A>0$ on $\R$). Then
$\Psi(x,u)=\Lambda(s(x)+iu)-\Lambda(s(x))-iux$, with the principal logarithm, satisfies
$G^M=e^{M\Psi}$ there, is real-analytic in $x$ and analytic in $u$, and
\begin{equation}\label{eq:Psi}
 \Psi=-\frac{v}{2}u^2-\frac{i\kappa_3}{6}u^3+\frac{\kappa_4}{24}u^4+\Psi_5,\qquad
 |\partial_x^j\Psi_5|\le C|u|^5,\qquad j\le2,\ |u|\le\varepsilon,
\end{equation}
because the $x$-derivatives of a Taylor remainder are the Taylor remainders of the
$x$-derivatives, which are bounded uniformly by Cauchy's estimates. Shrinking $\varepsilon$,
$\Re\Psi\le-cu^2$. By the cancellations of A5 (formalised), $|\partial_x\Psi|+
|\partial_x^2\Psi|\le Cu^2$. Put $u=z/\sqrt M$ and $Z=M^{1/12}$.

On $Z<|z|\le\varepsilon\sqrt M$,
$|\partial_x^je^{M\Psi}|\le C\{1+M(|\partial_x\Psi|+|\partial_x^2\Psi|)+M^2|\partial_x\Psi|^2\}
e^{M\Re\Psi}\le C(1+z^4)e^{-cz^2}$, whose integral is $O(e^{-cZ^2/2})=O(e^{-cM^{1/6}/2})$.

On $|z|\le Z$, write $M\Psi=-\frac v2z^2+w$ with
\begin{equation}\label{eq:w}
 w=-\frac{i\kappa_3z^3}{6\sqrt M}+\frac{\kappa_4z^4}{24M}+M\Psi_5\Bigl(x,\frac z{\sqrt M}\Bigr),
 \qquad |\partial_x^jw|\le C\,\frac{(1+|z|)^5}{\sqrt M}\le CM^{-1/12},\qquad j\le2 .
\end{equation}
Let $E=e^w-1-w-\frac{w^2}2$. By Taylor's formula for the exponential and the chain rule,
for $j\le2$, $|\partial_x^jE|\le C\sum|w|^{3-k}\,\Pi_k$, where $\Pi_k$ is a product of $k$
factors $|\partial_x^{i}w|$, $i\le2$; each term is at most $C(1+|z|)^{15}M^{-3/2}$. Expanding
$w+\frac{w^2}2$ and moving every term of order $M^{-3/2}$ into the remainder,
\begin{equation}\label{eq:exp}
 e^{M\Psi}=e^{-vz^2/2}\Bigl[1-\frac{i\kappa_3z^3}{6\sqrt M}
 +\frac1M\Bigl(\frac{\kappa_4z^4}{24}-\frac{\kappa_3^2z^6}{72}\Bigr)+\mathcal E_M\Bigr],
 \qquad |\partial_x^j\mathcal E_M|\le C\,\frac{(1+|z|)^{18}}{M^{3/2}} ,
\end{equation}
for $j\le2$ and $|z|\le Z$. The $x$-derivatives falling on $e^{-vz^2/2}$ produce polynomial
factors in $z$, so after integration over $|z|\le Z$ the remainder contributes
$O(M^{-3/2})$ in $C^2(J)$. The $z^3$ term and its $x$-derivatives are odd in $z$ and
integrate to zero over the symmetric range. Completing the remaining Gaussian moments to
$\R$ costs $O(e^{-cZ^2})$. With $\int z^4e^{-vz^2/2}dz=3\sqrt{2\pi}v^{-5/2}$ and
$\int z^6e^{-vz^2/2}dz=15\sqrt{2\pi}v^{-7/2}$ we obtain
\begin{equation}\label{eq:JM}
 J_M=\sqrt{\frac{2\pi}{Mv}}\,\bigl(1+R_M\bigr),\qquad
 R_M=\frac{q}{M}+O_{C^2(J)}\bigl(M^{-3/2}\bigr),\qquad
 q=\frac{\kappa_4}{8v^2}-\frac{5\kappa_3^2}{24v^3},
\end{equation}
where \eqref{eq:outer} and the middle region were used; both are exponentially small
relative to $M^{-1/2}$. Since $q$ is smooth on $J$, $\|R_M\|_{C^2(J)}\le C/M$.

\emph{Step 4.} By \eqref{eq:rep} and \eqref{eq:JM},
$\log g_M+MI-\frac12\log\{M/(2\pi v)\}=\log(1+R_M)$, and
$\|\log(1+R_M)\|_{C^2(J)}\le C\|R_M\|_{C^2(J)}$ once $\|R_M\|_{C^0}\le\frac12$. This
proves \eqref{eq:L1}, and indeed the sharper statement that the left side equals
$q/M+O(M^{-3/2})$ in $C^2(J)$.
\end{proof}

\begin{remark}
The sharper statement is confirmed numerically (Table~\ref{tab:saddle}). At the points
$x=0.05$, $0.15$, $0.25$, $0.35$ and $0.42$ one has $q=-0.151544$, $-0.163813$,
$-0.184391$, $-0.167178$ and $-0.089240$, while $M$ times the left side of \eqref{eq:L1}
at $M=800$ equals $-0.15157$, $-0.16386$, $-0.18449$, $-0.16724$ and $-0.08917$.
\end{remark}

\begin{remark}
For the theorem with an $o(1)$ error, only two consequences of Lemma~\ref{lem:L1} are used:
$\sup_{[0,b]}|\phi_M'-I'|\to0$ (with rate $O(M^{-1})$ for (9) of the note) and
$\sup_M\sup_{[-b,b]}|\phi_M''|<\infty$. An alternative route avoids differentiating under
the integral sign. The exact identities $f_M'(y)=f_{M-1}(y+\frac12)-f_{M-1}(y-\frac12)$ and
$f_M''(y)=f_{M-2}(y+1)-2f_{M-2}(y)+f_{M-2}(y-1)$ express $\phi_M'$ and $\phi_M''$ through
ratios of densities at shifted points. The undifferentiated case $j=0$ of the proof above,
applied to $f_{M-k}$ with $k\le2$, then suffices.
\end{remark}

\subsection{Lemma 2: corrected statement.}\label{sec:L2}

\begin{lemma}\label{lem:L2}
Let $H(a)=\frac1\pi\int_0^{\pi/2}a\sin\theta\,I'(a\sin\theta)\,d\theta$ and let $H_M$ be
defined in the same way with $\phi_M'$ in place of $I'$.
\begin{enumerate}
\item[(a)] $H$ is continuous and strictly increasing on $[0,\frac12)$, $H(0)=0$ and
$H(a)\to\infty$ as $a\uparrow\frac12$; hence $a_t$ is well defined for every $t>0$.
\item[(b)] For every $M\ge2$ and $t>0$ the equation $H_M(a)=t$ has a unique root
$a_{t,M}\in(0,\frac12)$; for $M=1$ it has none.
\item[(c)] For every compact $K\subset(0,\infty)$ there are $M_0,C$ with
$|a_{t,M}-a_t|\le C/M$ for all $t\in K$ and $M\ge M_0$.
\end{enumerate}
\end{lemma}

\begin{proof}
Part (a) is in the note, with the divergence argument of A9. For (b), $x\phi_M'(x)$ is
continuous on $[0,\frac12)$, vanishes at $0$ and is strictly increasing (A11 for $M\ge3$;
$x\phi_2'(x)=x/(1-2x)$ for $M=2$). Near $\frac12$ the density is a positive multiple of
$(\frac12-x)^{M-1}$, so $x\phi_M'(x)\ge c/(\frac12-x)$ there and $H_M(a)\to\infty$; the
root exists and is unique. For (c), choose $0<\mathrm{lo}<\mathrm{hi}<b<\frac12$ with
$H(\mathrm{lo})<\min K$ and $\max K<H(\mathrm{hi})$, and let $m=\min_{[0,b]}I''>0$. Lemma~1
gives $|\phi_M'-I'|\le C/M$ on $[0,b]$, and $I'\ge0$, $I'(y')-I'(y)\ge m(y'-y)$ on $[0,b]$.
The formal theorem \texttt{endpoint\_lemma} then yields (c) with constant $2bC/(m\,\mathrm{lo})$.
\end{proof}

\subsection{Lemma 3: corrected statement.}\label{sec:L3}

\begin{lemma}\label{lem:L3}
Let $K\subset(0,\infty)$ be compact. There are $N_0$ and constants $r^*>1$,
$\alpha\in(0,\frac14]$, $C>0$, $\alpha_1=\frac34$, $C_1>0$ and an open neighbourhood $I_0$ of
$\pm1$ such that, for all $(N,M)$ with $N\ge N_0$, $M\ge3$ and $N/M\in K$, the weight
$e^{-2NQ_{N,M}}=a_{N,M}g_M(a_{N,M}\,\cdot)$ on $I_{N,M}=(-1/(2a_{N,M}),1/(2a_{N,M}))$
satisfies conditions (i)--(v) of Definition~1.1 of \cite{levinlubinsky2018} with these
constants.
\end{lemma}

\begin{proof}
Let $a_\pm$ be the extreme values of $a_t$ on $K$ and fix $\delta>0$ and $r^*>1$ with
$r^*(a_++\delta)<\frac12$; by Lemma~\ref{lem:L2}, $a_-/2\le a_{N,M}\le a_++\delta$ for
$N\ge N_0$. Then $[-r^*,r^*]\subset I_{N,M}$ (formalised: \texttt{interval\_contains}).
Condition (i) holds by A13 and condition (ii) by A11. By \eqref{eq:Q},
$|Q'_{N,M}(x)-Q'_{N,M}(y)|\le\frac{a^2M}{2N}L_M|x-y|$ on $[-r^*,r^*]$, where
$L_M=\sup\{|\phi_M''(u)|:|u|\le\frac12-\eta\}$ for $\eta=\frac12-r^*(a_++\delta)$; $L_M$ is
bounded uniformly in $M\ge M_0$ by Lemma~\ref{lem:L1}, and $M/N\le1/\min K$. Hence $Q'_{N,M}$
is uniformly Lipschitz on $[-r^*,r^*]$, and therefore uniformly H\"older of orders
$\alpha$ and $\alpha_1$ there; take $I_0=(-1-\eta',-1+\eta')\cup(1-\eta',1+\eta')$ with
$1+\eta'<r^*$. These give (iii) and (iv). Condition (v) is \eqref{eq:mrs} together with
the second MRS equation, which holds by parity; by Remark~(b) after Definition~1.1 these are
equivalent to (v).
\end{proof}

\subsection{Proof of the theorem.}\label{sec:thm}

Let $E(N,M)=\beta_N^{[M]}/M^2-F(N/M)$, and let $K'=[t_1,t_2]\subset(0,\infty)$ with $t_1<t_2$
and $K\subset K'$. Take any $m:\N\to\N$ with $n/m(n)\in K'$ for all large $n$. Define
$Q_n=Q_{n,m(n)}$ for $n\ge N_0(K')$ and $Q_n(x)=x^2$ on $\R$ otherwise. By
Lemma~\ref{lem:L3}, $\{Q_n\}\in\Q$, so Theorem~13.4 of \cite{levinlubinsky2018} gives
$A_{n,n-1}=\frac12+o(1)$. By \eqref{eq:scaling}, $E(n,m(n))=a_{n,m(n)}^2A_{n,n-1}^2-a_{t}^2/4$
with $t=n/m(n)$. Since $a_{n,m(n)}-a_t\to0$ by Lemma~\ref{lem:L2}(c), and $F$ is bounded,
$E(n,m(n))\to0$. This is the hypothesis \texttt{hseq} of the formal theorem
\texttt{uniform\_of\_sequential}. The filling hypothesis is \texttt{fill\_ratio}, so that
theorem yields: for every $\varepsilon>0$ there is $L$ with $|E(N,M)|<\varepsilon$ whenever
$N\ge L$ and $N/M\in K$. This is (2) of the note.\qed

\subsection{What remains open.}\label{sec:open}

On paper, nothing remains open: with Sections~\ref{sec:L1}--\ref{sec:thm} the theorem is
proved modulo the published Theorem~13.4 and Remark~(b) after Definition~1.1 of
Levin--Lubinsky, which are used as black boxes. No rate is claimed. A rate
$O(N^{-\tau})$ would require knowing that the implied constants in Theorem~13.4 depend on
$\{Q_n\}$ only through the constants of Definition~1.1. The book does not state this
explicitly, and I have not verified it through the proofs.

\section{Search for counterexamples and numerical consistency checks.}\label{sec:num}

\textbf{Method.} The weight is a polynomial of degree $M-1$ on each knot interval. Gauss--Legendre
quadrature with $n_g\ge N_{\max}+M/2$ nodes per knot interval therefore reproduces the
first $N_{\max}$ recurrence coefficients exactly (the discretisation idea of Gautschi
\cite{gautschi2017}). Spline values come from the positive B-spline recurrence, and the
discretised Stieltjes procedure is run in 64-bit-mantissa extended precision. Against exact
rational values (moments and the Chebyshev algorithm) the code agrees to about 17
significant digits; for example both give $\beta_{12}^{[40]}=33.68595680113984\ldots$. All
diagonal coefficients vanish to within $10^{-16}$. The profile $F$ was computed with
\texttt{mpmath} at 30 digits. Scripts and logs are in the folder
\texttt{numerics/uniform\_lanczos} of the certification repository.

\begin{table}[h]
\centering
\begin{tabular}{@{}cccccc@{}}
\toprule
$t=N/M$ & $F(t)$ & $M\,E$, $M=100$ & $M\,E$, $M=200$ & $M\,E$, $M=400$ & $E$, $M=400$\\
\midrule
$1/4$ & $0.0178432$ & $0.010828$ & $0.010800$ & $0.010786$ & $2.70\times10^{-5}$\\
$1/2$ & $0.0303768$ & $0.017171$ & $0.017111$ & $0.017081$ & $4.27\times10^{-5}$\\
$1$   & $0.0445361$ & $0.017648$ & $0.017602$ & $0.017579$ & $4.39\times10^{-5}$\\
$2$   & $0.0548091$ & $0.010479$ & $0.010471$ & $0.010466$ & $2.62\times10^{-5}$\\
$4$   & $0.0598372$ & $0.004340$ & $0.004341$ & $0.004342$ & $1.09\times10^{-5}$\\
\bottomrule
\end{tabular}
\caption{$E=\beta_N^{[M]}/M^2-F(N/M)$ with $N=tM$ (largest case $N=1600$). $ME$ is
essentially constant in $M$, i.e.\ $E=c(t)/M+O(M^{-2})$.}\label{tab:beta}
\end{table}

\begin{table}[h]
\centering
\begin{tabular}{@{}ccccc@{}}
\toprule
$x$ & $q(x)$ & $MR$ at $M=800$ & $MR'$, $M=50\to800$ & $MR''$, $M=50\to800$\\
\midrule
$0.05$ & $-0.151544$ & $-0.15157$ & $-0.06315\to-0.06186$ & $-1.2664\to-1.2388$\\
$0.25$ & $-0.184391$ & $-0.18449$ & $-0.19104\to-0.18107$ & $2.2461\to2.2481$\\
$0.42$ & $-0.089240$ & $-0.08917$ & $0.50738\to0.56497$ & $-35.856\to-35.183$\\
\bottomrule
\end{tabular}
\caption{$R=\log g_M+MI-\frac12\log\{M/(2\pi v)\}$ and its derivatives, computed exactly
through $f_M'=\Delta f_{M-1}$ and $f_M''=\Delta^2f_{M-2}$; $MR$, $MR'$ and $MR''$ converge,
as Lemma~\ref{lem:L1} predicts, and $MR\to q$.}\label{tab:saddle}
\end{table}

\textbf{Counterexample search.} I looked for failures of uniformity near particular ratios
($t$ from $\frac14$ to $4$), for normalisation errors (a wrong factor would show up as a
non-vanishing limit of $E$, whereas $E\to0$ like $1/M$ in every case), and for
finite-parameter behaviour incompatible with the limit. I found none. The edge of $J$
($x=0.42$) has larger constants, as expected, but the same $1/M$ behaviour. These
computations are evidence, not proof.

\textbf{Milovanovi\'c's Conjecture~2.} As a by-product, for fixed $k=3,6,10$ the error of
the five-term formula of Conjecture~2 behaves like $c_k\,m^{-4}$ for $m=50$ to $400$ (for
$k=3$, $m^4\times\text{error}$ is $0.0100$, $0.0099$, $0.0099$, $0.0098$). This supports the
conjecture numerically and proves nothing.

\section{Originality and bibliography.}\label{sec:orig}

\textbf{(I) Fixed $M$, $N\to\infty$.} For fixed $M\ge2$ the weight $f_M$ is positive on
$(-\frac M2,\frac M2)$ and vanishes polynomially at the endpoints, so it lies in the Szeg\H{o}
class, and a fortiori Rakhmanov's theorem (see \cite{denisov2004}) applies. Hence
$\beta_N^{[M]}\to(M/4)^2=M^2/16$ as $N\to\infty$. This is classical and says nothing uniform
in $M$, as the note correctly states.

\textbf{(II) Milovanovi\'c's conjectures.} I read Section~5 of \cite{milovanovic2017} in the
original (pp.~345--347). Conjecture~1 asserts that for $k\ge3$,
$\beta_k^{(m)}=C_k\,q_{s_{k-3}}(m)q_{s_k}(m)/\{q_{s_{k-2}}(m)q_{s_{k-1}}(m)\}$ with integer
polynomials of degrees $s_k=\{2k^2-1+(-1)^k\}/8$ and rational $C_k$ (an exact-structure
statement at fixed $k$). Conjecture~2 asserts, for fixed $k\ge1$ and $m\to\infty$,
\begin{multline}\label{eq:conj2}
 \beta_k^{[m]}\asymp\frac{k}{12}m-\frac{k(k-1)}{20}+\frac{4k(k-1)(k-2)}{525m}
 +\frac{2k(k-1)(k-2)(5k-9)}{2625m^2}\\
 +\frac{2k(k-1)(k-2)(552k^2-1279k+255)}{1010625m^3}+\cdots .
\end{multline}
The theorem under review proves neither conjecture, and nothing here should be read as
claiming that it does. There is, however, an exact formal consistency, which I verified by
exact rational arithmetic. The Taylor coefficients of $F$ at $0$ are
\begin{equation}\label{eq:Fseries}
 F(\alpha)=\frac{\alpha}{12}-\frac{\alpha^2}{20}+\frac{4\alpha^3}{525}+\frac{2\alpha^4}{525}
 +\frac{368\alpha^5}{336875}-\frac{544\alpha^6}{1340625}+O(\alpha^7),
\end{equation}
and the first five coincide with the coefficients of the top-degree monomials
$k^jm^{2-j}$ in the five terms of \eqref{eq:conj2}; note that
$\frac{2\cdot5}{2625}=\frac{2}{525}$ and $\frac{2\cdot552}{1010625}=\frac{368}{336875}$. This
is what one expects if the fixed-$k$ expansion and the ratio-scaling limit match in an
overlap regime $1\ll k\ll m$. It is a non-trivial check of both, but it does not prove
either, because no uniformity bridging the two regimes is available. If Conjecture~2
continues with further terms, \eqref{eq:Fseries} predicts the top-degree coefficient
$-544/1340625$ of the next one.

\textbf{(III) The two-parameter theorem.} The general mechanism, that recurrence
coefficients are asymptotic to half the Mhaskar--Rakhmanov--Saff number, goes back to
Freud's conjecture for fixed exponential weights \cite{lms1988}. For varying weights the
tool is \cite{levinlubinsky2018}. The theorem is therefore a new \emph{application} of an
established theorem, not a new general theorem. Its specific content is: the identification
of the external field with the Cram\'er rate function $I$ of the uniform law; the uniform
$C^2$ saddlepoint control of the cardinal B-spline family (Lemma~\ref{lem:L1}, of the
classical type initiated by Daniels \cite{daniels1954}); and the explicit profile $F$. I
searched the web, Springer and arXiv, starting from \cite{milovanovic2017} and Gautschi's
companion paper \cite{gautschi2017}. I found no earlier statement of this joint limit or of
$F$. The abstract of \cite{gautschi2017} concerns computation only; its full text was not
accessible to me. A limited search cannot establish priority, so the result should be
described as ``to the author's knowledge new''.

\section{Formal-machine certification (Lean~4).}\label{sec:lean}
\sloppy

\textbf{Environment.} Lean \texttt{v4.35.0}, Mathlib commit
{\small\texttt{0272bbe1c093d038d7257ab15306abd58245f9ff}}\linebreak (the same as for the certification
repository of V478). Library \texttt{UniformLanczos} (ten modules) in the repository
\texttt{ralfanop/Certifications\_Eulerian\_polar-curve\_computations}, branch
\texttt{hjm2026-check}, folder \texttt{lean/UniformLanczos}. Every file compiles with no
error and no \texttt{sorry}, and no \texttt{native\_decide} is used. All 61 main
declarations depend only on Lean's standard axioms \texttt{propext},
\texttt{Classical.choice} and \texttt{Quot.sound} (file \texttt{lean/ULAxiomCheck.lean},
log \texttt{lean/logs/uniform\_lanczos\_axioms.log}). A clean build of the library takes
about 95~seconds.

\textbf{What is machine-checked.} The following list gives the module and the
main declarations for each part of the note.
\begin{enumerate}
\item \texttt{Cancellation}: the derivative cancellations of Lemma~1 (A5) and Legendre
duality $I'=s$, $I''=1/v$, for arbitrary $\Lambda$ with the stated derivatives
(\texttt{tilt\_first}, \texttt{tilt\_second}, \texttt{tilt\_cancellations},
\texttt{legendre\_deriv}, \texttt{legendre\_second}).
\item \texttt{Endpoint}: all of Lemma~2(c) and the monotonicity and continuity in
Lemma~2(a)--(b), in the $\theta$-form of $H$ (\texttt{Hint\_strictMonoOn},
\texttt{Hint\_continuousOn}, \texttt{Hint\_slope}, \texttt{Hint\_close},
\texttt{exists\_unique\_root}, \texttt{root\_stability}, \texttt{endpoint\_lemma}).
\item \texttt{ClassQ}: the elementary parts of Lemma~3, namely the formula for $Q'$, strict
monotonicity, Lipschitz transfer, Lipschitz~$\Rightarrow$~H\"older and
$[-r^*,r^*]\subset I_{N,M}$ (\texttt{Q\_hasDerivAt}, \texttt{Qprime\_strictMono},
\texttt{Qprime\_lipschitz}, \texttt{holder\_of\_lipschitz}, \texttt{interval\_contains}).
\item \texttt{MRS}: the equivalence of the first MRS equation (arcsine form) with the
endpoint equation, the second MRS equation by parity, and the scaling
$\beta=(Ma)^2A^2$ (\texttt{cos\_odd\_eq\_two\_sin}, \texttt{mrs\_first\_iff},
\texttt{mrs\_second}, \texttt{beta\_scaling}).
\item \texttt{Diagonal}: the uniformity argument of the theorem and the filling lemma
(\texttt{uniform\_of\_sequential}, \texttt{fill\_ratio}).
\item \texttt{Corollary}: $\int_0^{\pi/2}\sin^4=3\pi/16$, the expansion
$H(a)=3a^2+\frac{27}5a^4+O(a^6)$ from that of $I'$, and the inversion to
$F(\alpha)=\frac\alpha{12}-\frac{\alpha^2}{20}+O(\alpha^3)$
(\texttt{Hint\_expansion}, \texttt{A\_le\_half\_alpha}, \texttt{inversion}).
\item \texttt{Lambda} (added in the revision): the concrete $\Lambda'(s)=\frac12+\frac1{e^s-1}-\frac1s$ is odd,
an increasing homeomorphism $\R\to(-\frac12,\frac12)$, with $0<\Lambda''\le\frac14$,
$0<\Lambda'(s)\le\frac s4$ and $\frac12-\frac1s<\Lambda'(s)\le\frac12-\frac1{s+2}$ for $s>0$; its
inverse $I'$ is continuous and strictly increasing, with $I'(y')-I'(y)\ge4(y'-y)$ and
$I'(y)\ge\frac1{1/2-y}-2$.
\item \texttt{EndpointConcrete} (added): $H(a)\to\infty$ as $a\uparrow\frac12$, and Lemma~2(c) with every
property of the limit proved; only the bound $|\phi_M'-I'|\le C/M$ (Lemma~1) and the regularity of
$\phi_M'$ (Floater) remain hypotheses (\texttt{Hint\_unbounded}, \texttt{endpoint\_concrete}).
\item \texttt{Arcsine} (added): the substitution $y=\cos\theta$ in the improper integrals (1.13)--(1.14)
of Levin--Lubinsky, for Lebesgue integrals and with no integrability hypothesis, and the equivalence of
(1.13) with the endpoint equation (\texttt{arcsine\_substitution}, \texttt{mrs\_LL\_first\_iff},
\texttt{mrs\_LL\_second}).
\item \texttt{Expansion} (added): $|\Lambda'(s)-\frac s{12}+\frac{s^3}{720}|\le s^5/25600$ for $0<s\le2$,
$|I'(x)-12x-\frac{144}5x^3|\le520x^5$ for $0\le x\le\Lambda'(1)$, and unconditionally
$F(\alpha)=\frac\alpha{12}-\frac{\alpha^2}{20}+O(\alpha^3)$ (\texttt{Lp\_taylor}, \texttt{Ip\_expansion},
\texttt{F\_expansion}).
\end{enumerate}

\textbf{What is \emph{not} machine-checked.} The following inputs are hypotheses of the
formal statements, not formal theorems.
\begin{enumerate}
\item The analytic content of Lemma~1, i.e.\ the representation \eqref{eq:rep}, the bounds on
$G$ and the Laplace-method estimates. Consequently the hypotheses
$|\phi_M'-I'|\le C/M$ and $\sup|\phi_M''|<\infty$ are not formal.
\item The log-concavity of B-splines (Floater).
\item All potential theory, i.e.\ Remark~(b) after Definition~1.1 and Theorem~13.4 of
Levin--Lubinsky. Neither is in Mathlib, and formalising them is far beyond the scope of
this review.
\end{enumerate}
Hence the three lemmas are \emph{not} formalised in full. Lemma~2 is formalised completely
except for its two finite-$M$ inputs (the bound of Lemma~1 and Floater's log-concavity). Lemmas~1 and~3 are formalised only in their algebraic and
elementary parts. The deduction of the theorem from them is formalised as logic, with
Levin--Lubinsky's theorem entering as an explicit hypothesis.

\section{Addendum: revision of the same day.}\label{sec:addendum}
At the author's request every defect listed in Section~\ref{sec:audit} has been corrected in a
separate revised note, \emph{Uniform Two-Parameter Asymptotics for Lanczos Coefficients of Centered
Uniform Convolutions} (file \texttt{UL\_Note\_Revised.pdf}, TJM template of V035, V032 and V035
unchanged). It contains the complete proof of Lemma~1 of Section~\ref{sec:L1}, with the sharper
remainder $q/M+O(M^{-3/2})$. It restricts the endpoint and class-$\mathcal Q$ lemmas to $M\ge2$ and
$M\ge3$, corrects the choice of $r^*$, and makes the filling of the diagonal sequence explicit. It
records Levin--Lubinsky's recurrence convention and the misprint in Theorem~2.3(c), adds the
six-term expansion of $F$ with its consistency with Conjecture~2, and separates the result from
the fixed-degree Theorem~5.1 of V035. The formal coverage was extended as listed in
Section~\ref{sec:lean}: all properties of the limiting objects are now machine-checked, and the
remaining hypotheses concern only the finite-$M$ functions (Lemma~1 and Floater's inequality) and
the potential theory of Levin--Lubinsky.

\section{Final recommendation.}\label{sec:rec}

The theorem is, in my judgement, correct, and with the repairs of
Section~\ref{sec:proof} it has a complete proof on paper conditional on
\cite[Theorem~13.4]{levinlubinsky2018}. The note is not ready for incorporation as
it stands. Before incorporation, five obligations remain:
\begin{enumerate}
\item replace the proof of Lemma~1 by a complete one, such as Section~\ref{sec:L1};
\item restrict Lemma~2 to $M\ge2$, and Lemma~3 to $M\ge3$ and $N\ge N_0(K)$;
\item correct the choice of $r^*$ (A13);
\item make the filling of the diagonal sequence explicit (A16);
\item describe the result as a new application of Levin--Lubinsky's theorem, with
\cite{milovanovic2017} and \cite{gautschi2017} cited and the consistency
\eqref{eq:Fseries} with Conjecture~2 stated as consistency only.
\end{enumerate}
These five obligations are discharged in the revised note described in
Section~\ref{sec:addendum}. Separately, the author's own criterion was that the three lemmas enter
the manuscript only if they are formalised and certified. That criterion is \emph{not} met. Lemma~2
is machine-checked except for its two finite-$M$ inputs. Lemmas~1 and~3 are only partly
machine-checked, and their remaining parts depend on Fourier analysis and potential theory
that are not available in Mathlib. Whether to include the theorem with a paper proof and
an explicit statement of its formal coverage, or to wait for a fuller formalisation, is the
author's decision.

\begin{thebibliography}{9}
\bibitem{daniels1954}
H.~E. Daniels, Saddlepoint approximations in statistics, \emph{Ann.\ Math.\ Statist.}\
\textbf{25} (1954), 631--650.
\bibitem{denisov2004}
S.~A. Denisov, On Rakhmanov's theorem for Jacobi matrices, \emph{Proc.\ Amer.\ Math.\ Soc.}\
\textbf{132} (2004), 847--852.
\bibitem{floater2024}
M.~S. Floater, Log-concavity of B-splines, \emph{J.\ Approx.\ Theory} \textbf{300} (2024),
106042,\newline\mbox{\url{https://doi.org/10.1016/j.jat.2024.106042}}.
\bibitem{gautschi2017}
W.~Gautschi, Polynomials orthogonal with respect to cardinal B-spline weight functions,
\emph{Numer.\ Algorithms} \textbf{76} (2017), 1099--1107,\newline\mbox{\url{https://doi.org/10.1007/s11075-017-0298-9}}.
\bibitem{levinlubinsky2018}
E.~Levin and D.~S. Lubinsky, \emph{Bounds and Asymptotics for Orthogonal Polynomials for
Varying Weights}, SpringerBriefs in Mathematics, Springer, Cham, 2018,\newline\mbox{\url{https://doi.org/10.1007/978-3-319-72947-3}}.
\bibitem{lms1988}
D.~S. Lubinsky, H.~N. Mhaskar and E.~B. Saff, A proof of Freud's conjecture for exponential
weights, \emph{Constr.\ Approx.}\ \textbf{4} (1988), 65--83.
\bibitem{milovanovic2017}
G.~V. Milovanovi\'c, Symbolic--numeric computation of orthogonal polynomials and Gaussian
quadratures with respect to the cardinal B-spline, \emph{Numer.\ Algorithms} \textbf{76}
(2017), 333--347,\newline\mbox{\url{https://doi.org/10.1007/s11075-016-0256-y}}.
\end{thebibliography}

\end{document}
